\section{Object Detection và YOLO}

\subsection{Bài toán Object Detection}

Object Detection xác định đồng thời lớp và vị trí của các đối tượng trong ảnh.
Đầu ra thường là tập các bounding box $b_i=(x_i,y_i,w_i,h_i)$, nhãn lớp và độ
tin cậy. Khác với phân loại ảnh, một ảnh có thể chứa nhiều đối tượng và mô hình
phải giải quyết cả bài toán định vị lẫn phân loại.

Độ chồng lấp giữa hộp dự đoán $A$ và hộp chuẩn $B$ được đo bằng Intersection
over Union:
\begin{equation}
  IoU(A,B)=\frac{|A\cap B|}{|A\cup B|}.
\end{equation}
Một dự đoán được xem là đúng khi lớp phù hợp và IoU vượt ngưỡng đánh giá.

\subsection{Kiến trúc YOLO}

YOLO là họ mô hình phát hiện đối tượng một giai đoạn, thực hiện trích xuất đặc
trưng và dự đoán trong một đồ thị thống nhất \cite{redmon2016,redmon2017}.
Kiến trúc tổng quát gồm backbone để trích xuất đặc trưng, neck để tổng hợp đặc
trưng đa tỉ lệ và detection head để dự đoán vị trí, objectness và lớp.

Các biến thể YOLO khác nhau ở block tích chập, cơ chế kết hợp đặc trưng, hàm
loss và cách biểu diễn detection head. Khi tối ưu, cần bảo toàn những tầng ảnh
hưởng mạnh đến đầu ra, đồng thời ưu tiên cắt tỉa cấu trúc mà TensorRT có thể khai
thác thành giảm thời gian thực thi thực tế.

\subsection{Quy trình huấn luyện YOLO}

Quy trình huấn luyện gồm chuẩn hóa nhãn, chia train--validation--test, tăng cường
dữ liệu, khởi tạo trọng số, tối ưu loss và chọn checkpoint. Mọi thí nghiệm phải
lưu seed, phiên bản mã nguồn, cấu hình optimizer, learning rate, batch size, image
size và số epoch để có thể tái lập.

Tập test chỉ được dùng cho đánh giá cuối. Các quyết định về pruning ratio, lịch
fine-tuning và ngưỡng dừng sớm phải dựa trên tập validation nhằm tránh điều chỉnh
gián tiếp theo tập test.

\subsection{Các chỉ số đánh giá}

Với $TP$, $FP$ và $FN$ lần lượt là số dự đoán đúng, dương tính giả và âm tính
giả, Precision và Recall được xác định bởi:
\begin{equation}
  Precision=\frac{TP}{TP+FP}, \qquad
  Recall=\frac{TP}{TP+FN}.
\end{equation}

Average Precision là diện tích dưới đường Precision--Recall của một lớp. mAP
là trung bình AP trên các lớp. $mAP@50$ dùng ngưỡng IoU 0,5; $mAP@50{:}95$
lấy trung bình trên các ngưỡng từ 0,5 đến 0,95 với bước 0,05. Đồ án báo cáo đồng
thời Precision, Recall, $mAP@50$ và $mAP@50{:}95$ để tránh kết luận dựa trên
một chỉ số đơn lẻ.